\documentclass[11pt]{article}
\usepackage[utf8]{inputenc}
\usepackage[T1]{fontenc}
\usepackage{lmodern}
\usepackage[margin=1in]{geometry}
\usepackage{csquotes}
\usepackage[style=authoryear,backend=biber,sorting=nyt,maxbibnames=99,giveninits=true,uniquename=init]{biblatex}
\usepackage[hidelinks]{hyperref}
\addbibresource{references.bib}
\begin{document}
\section*{Data and Experimental Validation Needs}

\subsection{Overview}

A recurring conclusion across the tile-drainage literature is that theory,
models, and management hypotheses have outrun the field data available to
validate them. The studies surveyed here converge on a shared research
frontier defined less by a lack of ideas than by a lack of evidence: process
representations that remain untested outside the conditions in which they were
built, models whose parameters cannot be constrained without measurements that
were never made, and management practices whose benefits are asserted from
short records at single sites. This section synthesises the identified data and
experimental-validation needs. Several broad lines of evidence emerge. First,
process-based drainage and water-quality models (DRAINMOD and its derivatives,
SWAT, APEX, HYDRUS, CRHM, and numerous salinity codes) are repeatedly flagged as
requiring further testing across soils, climates, crops, and management systems,
with particular weaknesses in phosphorus, macropore/preferential-flow, and
seepage components \parencite{Youssef2006,Askar2021b,Askar2025,Radcliffe2015}.
Second, there is a strong and near-universal call for long-term, year-round, and
multi-site monitoring to close water and nutrient budgets
\parencite{Kladivko2001,Sharpley2019,Frankenberger2024}. Third, rapidly maturing
remote-sensing and geophysical methods for mapping drainage networks need
systematic evaluation of the site conditions under which they succeed
\parencite{Carlsen2024,Wienken2024}. Fourth, emerging management
strategies -- controlled drainage, drainage water recycling, and edge-of-field
treatment -- lack the coordinated field-scale experiments and simulations needed
to quantify their trade-offs \parencite{Hay2021,Youssef2018,Mendes2020}. Finally,
authors repeatedly emphasise improvements to experimental design, statistical
methodology, uncertainty quantification, and the reporting of site
characterisation, management, and economic data
\parencite{Pagliari2021,Singh2022b,Gates2002}. The subsections below organise
these needs by domain.

\subsection{Validation and improvement of process-based drainage models}

The most frequently cited need is more rigorous, independent testing of
simulation models beyond the conditions in which they were calibrated.
\textcite{Youssef2006} note that further research is needed to test the model for
different management practices and soil and climatological conditions, and that a
more rigorous field testing should be conducted to validate application to
drained lands -- a need they describe as common to nearly all simulation models.
Similar generic calls for further validation, broader-condition testing, and
parameter sensitivity analysis recur for pesticide, salt, and flow models
\parencite{Branger2009,Hirekhan2007,Yang2008b,Xian2017}. \textcite{Yang2008b}
stress that soil-related coefficients -- particularly the drainage-volume/water-table-depth
relationship -- must be adjusted for local soil and water conditions before the
improved DRAINMOD extensions can be applied elsewhere, and propose integrating
remote sensing and GIS for input preparation and eventual incorporation into
decision-support tools. \textcite{Xian2017} likewise call for testing across
different climates, soils, and agronomic practices and for alternate approaches
to improve simulation of drainage peaks. \textcite{Mendez2004} argue that more
experimental studies in different soils and climates are needed to verify
modelling findings, and that long-term simulations of 50 to 100 years are
required to assess whether a given drainage design meets chosen criteria.

\subsubsection{DRAINMOD, controlled drainage, and drainage water recycling components}

Extensions of DRAINMOD are singled out for continued development and testing.
\textcite{Askar2021b} call for more testing of newly incorporated hydrological
and biogeochemical components -- especially the macropore flow, phosphorus cycle,
and sediment components -- with rigorous, quantitative assessment of performance
under varying water management, farming practices, and climatic conditions.
\textcite{Askar2025} extend this to the phosphorus component specifically, urging
validation across a range of soils, drainage characteristics, and climates,
recommending that key coefficients become user inputs for flexibility, and
highlighting the soil-erosion and macropore-flow components as priorities for
rigorous testing. For drainage water management and recycling (DWR/DWM),
\textcite{Moursi2022} call for testing DRAINMOD-DWR against field-measured data
and for developing routines that simulate the fate and transport of nitrogen,
phosphorus, and sediment in the storage reservoir; \textcite{Moursi2023} add that
the physical, chemical, and biological processes in the reservoir must be
investigated to optimise treatment efficiency and to identify conditions under
which reservoirs become nutrient sources, evaluated under varying site conditions
over longer periods. Related model-development needs include a time-variable
surface-depression-capacity input and expansion to wastewater and urban modules
\parencite{Abduljaleel2023}, and evaluation of a DRAINMOD-ANN model against
experimental data across different soil textures and field conditions
\parencite{Awad2022}. \textcite{Youssef2018} outline the experimental verification
required for controlled drainage: side-by-side comparison of controlled and free
drainage with extensive measurement of mineral-N concentrations in the root zone,
shallow groundwater at multiple depths, and drainage outflow, sustained long
enough to capture responses to precipitation, cropping, and fertilisation, plus a
separate simulation study to assess effects on vertical seepage across soils.
\textcite{Skaggs2010} emphasise that controlled-drainage experiments on drainage
volumes and N losses must be conducted at field or watershed scale so that
seepage impacts are properly represented, while \textcite{Skaggs2006} note that
additional research is needed to test relationships and develop equations for a
wider range of drain depths and surface storages before simplified drainage-design
methods can be accepted for general use. Complementary evaluation methods and
alternative Ks-measurement relationships are recommended alongside DRAINMOD in
data-scarce settings \parencite{NabiSichani2012}.

\subsubsection{Phosphorus, sediment, and macropore transport in water-quality models}

Phosphorus and its transport through macropores emerge as the least mature
component across water-quality models. \textcite{Radcliffe2015} provide the most
comprehensive agenda: field experiments using a nested, paired research design
are needed to improve phosphorus models for artificially drained fields, and
several codes require structural additions -- SWAT and ADAPT should route
macropore flow directly to tile drains or lateral flow to ditches, APEX needs a
macropore component, HSPF and P-indices in artificially drained areas should
simulate water and phosphorus movement to ditches or drains, DRAINMOD should be
expanded to simulate P losses to drains and P attached to sediment in runoff,
PLEASE needs a runoff component, and HYDRUS could be used experimentally to
understand transport to drains and parameterise simpler models -- with all model
studies quantifying the accuracy and uncertainty of predicted P losses.
Consistent with this, \textcite{Guo2020b} recommend modifying tile-drainage,
macropore-flow, and P-transport processes in the APEX/NTT interface, incorporating
physically based tile-drainage subroutines and macropore nutrient-transport
pathways, including particulate-P losses through tiles, and accounting for soil-test-P
effects over the full simulation period. \textcite{Ford2017} call for further
testing under varying fertiliser management (broadcast and incorporated, organic
and inorganic), additional field data to constrain subsurface-P parameter ranges
for tools such as the Nutrient Tracking Tool, and a better understanding of how
macropore flow and P loads respond to tillage, cover crops, fertiliser management,
and drainage water management. \textcite{BinShokrana2023b} identify specific
improvements to a DRAINMOD P component -- better partitioning of fertiliser P into
pools (especially under ponded conditions), improved macropore-flow simulation,
testing across climates, soils, crop varieties, and management, and refined
P-only performance statistics. \textcite{Gowda2012} anticipate further developing
the DRAINMOD suite to simulate phosphorus fate and transport in drained land.
Supporting empirical work is also required: \textcite{McDowell2001b} call for
further testing of CaCl$_2$-P as an estimator for dissolved reactive phosphorus in
drainage water from other soils and for determining how water and CaCl$_2$
extractions behave under different rainfall intensities and durations in the
field; \textcite{Morrison2013} recommend additional field research on how drainage
parameters affect runoff, subsurface drainage, and P losses, and improving
regression models with P saturation, soil-test P, and fertilisation rates;
\textcite{Sanford2016} require field-scale work with fluctuating flow rates and P
concentrations and tests using actual drainage effluent; and \textcite{Jia2012}
call for long-term monitoring of water-quality trends and investigation of the
sources and transport of P.

\subsubsection{Salinity, salt balance, and unsaturated-zone coupling}

In arid and irrigated settings, needs cluster around salt-transport validation and
coupling saturated and unsaturated processes. \textcite{Ebrahimian2015} recommend
applying models capable of simulating salt and nitrate concentrations in drainage
water and the soil profile to mitigate environmental impacts, while
\textcite{Sarangi2006} note that neural-network models can predict effluent
salinity and reclamation period from minimal input data but that SALTMOD may be
preferred for detailed salt-balance estimation. \textcite{Hirekhan2007} caution
that more evaluation studies are needed before firm recommendations on salt-transport
simulation can be made. Several Chinese studies press for tighter coupling of
water and salt dynamics: \textcite{Han2025} propose coupling future models with
unsaturated water-salt models such as HYDRUS-1D and UBMOD-S for more accurate
water and salt content in both zones and for developing regional subsurface-pipe
layout parameters; \textcite{Wang2025c} plan to integrate process-based models
(HYDRUS-2D/SWAP) with multi-objective optimisation and to consider climatic
variability, soil texture, groundwater mineralisation, and crop water stress; and
\textcite{Wang2024c} call for continuous model refinement, long-term simulation,
and exploration of additional factors affecting salt flux under varying
conditions. \textcite{Liu2022} recommend adding a detailed root-growth module,
automatic TDR monitoring of water and salt, evaluation under different irrigation
methods and salt levels, and validation with large field datasets, while
\textcite{Liu2026} extend simulation horizons to include climate-change scenarios,
long-term salt accumulation in deep soil and groundwater, and real-world
scalability. \textcite{Guo2025} similarly seek longer observation periods, a wider
range of environmental variables, optimal leaching quotas under diverse
conditions, and cost-benefit analyses, and further elucidation of the links among
surface salt accumulation, temperature, and groundwater depth. \textcite{Kelleners2000}
warn that model assumptions must be re-verified when applied to new areas and that
proper incorporation of regional seepage would require a fundamentally different
modelling approach. Earlier salinity-modelling series similarly staged validation
and salt-flow analysis across successive papers \parencite{KARAJEH1994}.

\subsubsection{Hydrological structure, seepage, and preferential flow representation}

A cross-cutting structural gap is the representation of seepage, groundwater
exchange, and preferential flow. \textcite{DaneshYazdi2016} suggest that
groundwater contributions could be incorporated through a double-storage model and
residual storage added if hydrochemical data were available, requiring further
data to refine interpretations. \textcite{Feng2018} argue that lateral seepage
must be allowed to vary during simulation and be included among model inputs, and
\textcite{Kelleners2000} echo the difficulty of representing regional seepage.
\textcite{Kompanizare2024} recommend adapting new tile-drainage modules to
different soil textures such as clay loam where preferential flow strongly affects
saturated storage and tile flow, adding explicit unsaturated-flow representation
where groundwater is disconnected from surface water, integrating with frozen-soil
and water-quality modules, and expanding to larger areas. \textcite{Reid2012}
stress that more work is needed to quantify the range of macropore flows to tile
drains across soil types and management systems, that seasonal (growing versus
non-growing) hydrologic regimes should be treated as discrete, and that a
framework is offered to encourage others to gather suitable data for testing.
\textcite{Allaire2011} call for testing each preferential-flow component
individually, validating them in the field, and analysing the individual impact of
each component on water quality, including extreme values. \textcite{Ojeda2018}
require a more accurate, independent method for determining daily subsurface
drainage from one layer to the next given preferential flow, plus sensitivity
analysis to identify key parameters needing measurement. \textcite{Ulén2016} note
that quantification of macropore flow paths is necessary for proper understanding
of subsurface-drain leaching losses and that a dynamic model accounting for frozen
soil and air temperature, together with measured data on soil structure and
soil-P, is urgently needed. \textcite{Frey2012} call for characterising near-surface
preferential flow paths across different soil, land-use, and climatic settings, and
\textcite{LESAFFRE1988} note that a better understanding of by-pass flow through
cracks is necessary to extend model application to full-year operation, alongside
continued field experimentation and automatic calibration of water-table shape
coefficients. Nitrogen-cycle and hydrological process representations attract
parallel needs: \textcite{Malone2017} suggest developing a surface-residue
component, a more detailed drainage approach, more thorough investigation of soil
field capacity for denitrification and N losses, and mobile/immobile water phases
to improve nitrate transport; \textcite{Gillette2018} call for more empirical data
on interactions between soil-water retention and the soil-N budget, including
mineralisation and denitrification, and for well-tested models and the use of
soft data and complete N budgets during calibration. Analytical and theoretical
advances -- rigorous theoretical analysis, formulas for sloped and multilayer
soils, and field examination -- are sought for improved subsurface-drainage designs
\parencite{Tao2016,Tao2019}, and time-step selection is flagged as a
study-specific choice requiring pre-assessment \parencite{Awad2022b}. Other
process needs include temporal variability in decay coefficients and deterministic
approaches for land-use effects on tile-drained water quality \parencite{Rudra2005},
improved transient and inversely calibrated representations \parencite{Sedaghatdoost2018},
and machine-learning surrogate models that expand training datasets across diverse
climatic and hydrogeological conditions, incorporate transient processes and
robust/stochastic optimisation, and integrate with real-time control
\parencite{Lei2026}.

\subsection{Long-term, year-round, and multi-scale monitoring}

Across nutrient, salt, and hydrological studies, authors emphasise that short,
seasonally incomplete records are inadequate. \textcite{Kladivko2001} argue that
year-round measurement of water flow and chemical concentration is necessary, for
at least several years on new sites, to determine which parts of the year should
receive management priority, and that tile flow versus surface runoff contributions
must be assessed at watershed scale. \textcite{Sharpley2019} present a broad agenda
calling for a more holistic understanding of subsurface hydrology, year-round data
collection to close water and nutrient balances, continued research on P movement
through drainage systems across soils in the Corn Belt, study of legacy N and P
pools and historical fertiliser effects, quantification of snowmelt-driven losses,
and evaluation of single, stacked, and edge-of-field practices and their
trade-offs. \textcite{Frankenberger2024} provide detailed monitoring
recommendations for drainage water management studies: high-frequency discharge,
surface-runoff, and seepage monitoring to capture the complete water budget;
process-based measurements including microbial activity, denitrification,
metabolism, soil-P pools, organic matter, infiltration rate, and redox potential;
consistent documentation of land management and DWM design (tile spacing, outlet
depth, timing of water retention); and a publicly available dataset supporting
re-analysis. \textcite{Williams2017} call for quantifying the relationship between
field-scale processes and watershed response, detailed management and soil-carbon
information as a priority in watershed water-quality studies, and further work on
dissolved organic carbon dynamics and sources including groundwater seepage and
in-stream production. \textcite{Drake2018} recommend field measurements of
discharge to validate hydrologic models and water monitoring over the entire year
even at reduced frequency, with continued monitoring and modelling to narrow
uncertainty in conservation-practice performance. Long-term monitoring is likewise
sought for edge-of-field and treatment systems -- accumulation and biofouling of
recycled-concrete-aggregate drains \parencite{McCulloch2017}, long-term maturation
of surface-flow constructed wetlands for P retention with attention to locally
available filter materials and cost-efficiency \parencite{Mendes2020}, long-term
performance of stormwater-retention devices \parencite{Xu2022}, and long-term and
negative groundwater-level effects with accumulated multi-field discharge data
\parencite{Okuda2020}. \textcite{Furukawa2008} call for careful evaluation of
leaching components and total carbon dynamics in relation to drainage systems,
including a quantification method for reduced carbon. \textcite{Frankl2023}
recommend tracking soil-property metrics as standard practice, controlling for
tillage, deeper sampling to resolve profile carbon changes, and pairing
water-quality studies with monitoring of changing physical and biological
properties over time.

\subsection{Transferability across soils, climates, crops, and management}

A pervasive limitation is that findings from single sites or narrow conditions
cannot yet be generalised. \textcite{Singh2022b} recommend that controlled-drainage/subirrigation
studies be conducted outside the U.S. Midwest and Canada across diverse climates,
crops, and soils; that additional studies compare irrigation methods in terms of
water supply and water-balance components; and that site characterisation --
including field slope, soil particle-size distribution by horizon, hydraulic
conductivity, drainage class, and soil moisture -- be reported, together with
measures of variability. \textcite{Bilotta2008} stress that work must be repeated
at larger scales and on different soil types as key modulating factors, and
\textcite{Goeken2015} call for investigations across geographic regions, diverse
perennial species and mixtures, and varied precipitation and weather.
\textcite{Nuutinen2001} recommend repeating earthworm-drainage studies on soils of
different texture and under less rainy conditions and including other earthworm
species, while \textcite{Nuutinen2011} seek horizontal-transect observations to
strengthen source-area interpretation and long-term relation of colonisation to
soil structure and hydrology. Climate-focused transferability needs include
high-resolution climate-change scenarios and effects on lodging, erosion, and
agrochemical transport \parencite{Rasmussen2018}; analysis of hydrological-model
contributions to total uncertainty using multiple models and assessment at
additional French sites under different climates and soils \parencite{Jeantet2022};
and validation of drainage-warming and soil-texture effects on crops via further
study and modelling \parencite{Jin2008}. Crop- and practice-specific extensions
are sought for paddy and trench drainage \parencite{Rahimi2024,Jia2023},
mole-drainage durability over longer periods and rainy seasons \parencite{JHA1995},
controlled-drainage principles suited to local crops and agronomy in new regions
\parencite{Sojka2020}, and interactions of crop rotation and tillage on subsurface
drainage \parencite{Ojeda2018}. \textcite{Yu2018} call for further field
evaluation of a salt-flushing method, comparison with intermittent ponding,
sprinkling, and drip irrigation, an analytical solution for generalisation, and
tests in complex heterogeneous soils. Assumptions must be re-verified whenever
models are transferred to new areas \parencite{Kelleners2000,Youssef2006}, and
localised design methods must be developed to reflect local agro-technical
conditions \parencite{Tlapáková2017}.

\subsection{Remote sensing and geophysical mapping of drainage networks}

Detecting and mapping buried drainage systems is a fast-moving frontier where
methods work in some conditions but lack systematic guidelines. \textcite{Wienken2024}
describe a functioning non-invasive detection method as indispensable, call for
urgent development of holistic and efficient methodologies for aging drainage
systems, and recommend combining UAV flights with ground-penetrating radar (GPR)
for improved mapping and characterisation. \textcite{Allred2020} lay out an
extensive UAV agenda: additional visible, multispectral, and thermal-infrared
surveys at multiple farm sites visited repeatedly over two years to define the
range of site conditions under which each imagery type maps drain lines; soil
temperature and emissivity measurements in small test plots of variable soil,
wetness, and cover; advanced processing (band stacking with decorrelation stretch,
minimum-noise-fraction and principal-component transforms); supervised
machine-learning algorithms to distinguish drain-line responses from field-operation
features; and decision-tree guidelines for practitioners. Related UAV/thermal work
seeks site-condition guidelines, test-plot quantification of surface thermal
behaviour, advanced image processing, and machine learning to differentiate
responses \parencite{Allred2021}, and extends thermal-infrared methods
to drainage-outlet detection under Midwestern conditions while assessing the
hydrologic impacts of outlets on ditches and streams \parencite{Allred2022}.
\textcite{Carlsen2024} recommend reconsidering the often-cited 2.54~cm minimum
rainfall for imagery acquisition, investigating deep learning and multitemporal
acquisition for drained-area segmentation and tile detection, exploring synthetic
aperture radar, determining the spatial resolution necessary for tile delineation,
performing more large-scale deep-learning studies, bridging computer-science and
geoscience practices through uniform evaluation standards, and expanding publicly
available drainage-map datasets. Machine-learning mapping across regions is further
advanced by calls for deep-learning and transformer-based encoder-decoder
architectures, higher-resolution imagery (harmonised Landsat/Sentinel-2, Planet,
or aerial), broader spatial and temporal application, and integration into
hydrologic and water-quality modelling frameworks \parencite{Wan2024}, and by
reducing tracing uncertainty through multiple interpreters, ground-truthed
datasets, additional raster covariates (soil type, surface cover, time of year),
wind-turbine training examples, and satellite time series \parencite{Redoloza2023}.
Geophysical evaluations echo these needs: near-surface geophysical data from more
sites and diverse soils \parencite{Allred2005}; systematic testing of antenna
configuration, frequency, and antenna-to-pipe orientation across soil types for
GPR \parencite{Allred2013}; prediction of GPR suitability via electromagnetic
simulation (e.g., gprMax) from EMI measurements and understanding dependence on
soil type, residue, tillage, ground wetness, and prior rainfall
\parencite{Koganti2020,Koganti2021}; concurrent EMI/GPR and complementary
drone-imagery and magnetic-gradiometer surveys \parencite{Koganti2020}; and
resistivity/electromagnetic tracer approaches designed for slope and known drain
orientation \parencite{Jouen2018}. Optical, radar, and topographic covariate
studies also seek more frequent imaging \parencite{Tlapáková2017}, repeated SAR
acquisitions with ground-based measurements to track crop growth across drainage
ridges \parencite{Meunier2001}, UAV-based cover-crop biomass and N-uptake mapping
across landscapes \parencite{Xu2018}, and richer topographic and soil covariates
in statistical drainage models \parencite{Harris2018}.

\subsection{Water balance, seepage, and groundwater interactions}

Closing the water balance -- particularly seepage and groundwater exchange --
remains a persistent data gap. \textcite{Grose2022} call for further exploration
to identify the contribution of groundwater relative to other sources during
stormflows, and \textcite{DaneshYazdi2016} would incorporate groundwater through
double-storage modelling. \textcite{Sedaghatdoost2018} note that providing
groundwater-table data in objective functions would help inverse modelling recover
soil properties more precisely, and \textcite{Feng2018} advocate incorporating
variable lateral seepage. Regional groundwater assessments call for further
radiocarbon sampling and continued rigorous evaluation of changing resource origins
\parencite{Tweed2018}, monitoring of groundwater chemistry at recharge zones of
well-fields for both shallow and deep flow systems \parencite{Dragon2016}, and
additional geologic and hydrologic data plus analysis of deep percolation through
silts and clays \parencite{BOOKMAN1959}. Aquifer characterisation needs include
accurate descriptions of the hydraulic properties of different aquifer layers
\parencite{KAMRA1991} and intensive investigation of the impervious-layer depth to
assess effluent-salinity variation under alternative drainage designs
\parencite{KAMRA1994}. \textcite{Stuetz2004} recommend greater use of test wells to
establish hydrochemical conditions and greater attention to water quality and well
design to minimise biofouling. \textcite{Schilling2022} suggest transient
simulations to test steady-state assumptions and further research on the capacity
of waterways to process excess water, and \textcite{Xu2022} note that updated local
groundwater-dynamics information is necessary for siting retention devices.
Watershed-scale salt and water studies plan continued data collection with
unsteady flow and mass-transport models, sensitivity and stochastic analysis, and
consideration of canal lining, horizontal subsurface drainage, and river-level
management \parencite{Gates2002,Burkhalter2006}. Groundwater as an alternative
water source and the effects of drainage from saline-alkali lands on groundwater
are flagged for investigation \parencite{Zheng2025b}, as are the interrelation of
irrigation practice and drainage design \parencite{Sinai2005} and long-term
groundwater-level effects of new drainage \parencite{Okuda2020}. \textcite{Keller2008}
suggest that long-term amelioration of nitrogen contamination may require perennial
deep-rooted crops and are investigating applicability at larger watershed scales.

\subsection{Experimental design, statistics, and uncertainty quantification}

Several authors identify methodological rather than substantive gaps.
\textcite{Pagliari2021} call for confirming whether beta distributions or time
series are the best approaches for analysing field-scale, unreplicated paired
drainage experiments, or whether new statistical methodologies are needed for
controlled-drainage studies. \textcite{Andrus2014} argue that directed
investigations of specific hypotheses are more productive than replicate spatial
analyses of additional biofilters and recommend more frequent shallow sampling at
specified spacings. \textcite{Okuda2020} note that field experiments with
replication are needed to enable statistical analysis, and \textcite{Wamono2016b}
recommend matching tractor power to implement size, comparing on a percentage
basis, and adding measurements of draft, slip, and trafficability. Uncertainty
quantification recurs strongly: stochastic and sensitivity analyses to address
parametric uncertainty \parencite{Gates2002,Burkhalter2006}, robust/stochastic
optimisation for parameter and input uncertainty \parencite{Lei2026}, explicit
estimation of prediction accuracy and uncertainty in phosphorus models
\parencite{Radcliffe2015}, reduction of interpreter-based mapping uncertainty
\parencite{Redoloza2023}, and analysis of the hydrological model's contribution to
total uncertainty using multiple models \parencite{Jeantet2022}. Improved sampling
and measurement designs are recommended for buffer and riparian studies, including
more sampling wells to capture spatial variability, tracer studies to attribute
flow proportions, and both surface and dissolved N$_2$O measurements
\parencite{Davis2019}, and dye tracing, control wells, and paired free-flowing
watershed comparisons at saturated-buffer sites \parencite{Jacquemin2020}. Improved
model performance is also sought through better evapotranspiration estimation and
additional predictor variables \parencite{Morrison2013}, additional SWAT performance
metrics, parameter refinement, longer observation periods, and larger datasets
\parencite{Zheng2025b}, and independent evaluation of calibrated models against
withheld data \parencite{Schoups2005}. \textcite{Schoups2005} further suggest
spatially distributing parameters to remove trade-offs between shallow and deep
water tables and multi-objective optimisation with additional criteria.

\subsection{Economic feasibility, management data, and databases}

A final cluster of needs concerns the socio-economic and data-infrastructure
context of drainage research. Economic feasibility analyses are repeatedly
requested: cost-benefit analysis of surface amendments and energy savings
\parencite{Wamono2016b}, regional economic benefits of drainage water recycling
under future climate and droughts with optimised reservoir dimensions
\parencite{Kosari2024}, overall economic benefit and inter-BMP interactions of
integrated systems \parencite{Li2021b}, comprehensive cost-benefit analysis of
subsurface drainage accounting for short- and long-term effects
\parencite{Guo2025}, economic feasibility of drainage as a preventive measure
against salinisation \parencite{Datta2000,HalbacCotoaraZamfir2022}, cost and
consumption in paddy fields \parencite{Jia2023}, cost-efficiency of edge-of-field
technologies \parencite{Mendes2020}, and cost reduction and technology
sophistication \parencite{Okuda2020}. Detailed management and soil data are
emphasised as prerequisites for reliable inference: verification of soil
information by field sampling and farmer records with accurate management data
\parencite{JORJANI1991}, detailed agricultural-management information and soil-carbon
change as watershed-study priorities \parencite{Williams2017}, sound pedological
and geological knowledge for interpreting drainage concentrations
\parencite{Jacobsen2007}, and richer site characterisation and variability
reporting \parencite{Singh2022b}. Data-infrastructure and regionalisation needs are
prominent in the German mapping work of \textcite{Hirt2005,Hirt2005b,Hirt2011},
which calls for comprehensive databases of tile-drained land proportions by soil
type, validation of existing databases, extrapolation from eastern to western
Germany and to European soil maps at the mesoscale, and inclusion of drained areas
in water- and nutrient-flux models so that mitigation measures can be designed; the
approach is transferable where representative drainage data and medium-scale soil
maps exist. Coordinated, coupled field-and-simulation research programmes are
called for to advance drainage water recycling across hydrology, crop-yield,
water-quality, complementary, and implementation categories, with monitoring across
sites and climates to document downstream load reductions and reservoir behaviour
\parencite{Hay2021}, alongside soil-moisture-sensor and remote monitoring of
control structures, validation under increasing climate variability, and
additional well-characterised sites \parencite{Willison2021}. Continued monitoring
of salinity and multi-scale water quality is essential to refine drainage-reuse
strategy \parencite{ElAgha2020}, and forestry, land-use, and topographic effects on
runoff generation and nitrogen retention require further analysis in Nordic and
Baltic settings \parencite{Deelstra2014,Deelstra2014b}. Site-specific reservoir and
constructed-wetland sizing \parencite{Haverstock2017}, coupling of soil-water
simulation with management-scenario testing \parencite{Bulot2017}, verification of
sulphur-loss and other process models under a range of climatic, management, and
soil conditions \parencite{SMITH1983}, and evaluation of forest water-management
systems and best management practices \parencite{MCCARTHY1992} round out the
distributed set of data and validation needs that define this research frontier.

\printbibliography
\end{document}
